\begin{theorem}[g-BQO]
For every nonidentity $g\in\mathsf{IncSeq}$ and every binary relation $R$
on a discrete type $A$,
\[
(A,R)\text{ is a $g$-better-relation}
\quad\Longleftrightarrow\quad
(A,R)\text{ is a better-relation}.
\]
In particular, if $Q$ is a quasi-order with relation $\leq$, then
\[
Q\text{ is $g$-bqo}\quad\Longleftrightarrow\quad Q\text{ is bqo}.
\]
\end{theorem}
\begin{proof}
Fix $g\neq\mathsf{IncSeqId}$.  By \noderef{MainProp},
\noderef{ContMorEq}, and the definition of \noderef{ContMor}, choose
continuous homomorphisms
\[
 K:(\mathsf{IncSeq},\mathsf{RightComp}(g))\longrightarrow
       (\mathsf{IncSeq},\mathsf{ShiftMap}),
 \qquad
 L:(\mathsf{IncSeq},\mathsf{ShiftMap})\longrightarrow
       (\mathsf{IncSeq},\mathsf{RightComp}(g)).
\]
Thus, by the two fields of \noderef{ContinuousHom}, $K$ and $L$ are
prefix-continuous and satisfy
\[
 K(\mathsf{RightComp}(g)(x))=\mathsf{ShiftMap}(K(x)),
 \qquad
 L(\mathsf{ShiftMap}(x))=\mathsf{RightComp}(g)(L(x)).
\]

We spell out how these maps transport complement homomorphisms.  Let
$H:(\mathsf{IncSeq},\mathsf{ShiftMap})\to(A,R^{\complement})$ be a
continuous homomorphism in the sense of \noderef{ContinuousRelHom}.  Define
$J(x)=H(K(x))$.  To check local constancy at $x$, first use local constancy
of $H$ to choose an output length $m$ such that every $z$ agreeing with
$K(x)$ through $m$ has $H(z)=H(K(x))$.  Prefix-continuity of $K$ then gives
an input length $k$ such that agreement with $x$ through $k$ forces
$K(z)$ to agree with $K(x)$ through $m$.  Consequently $J(z)=J(x)$ for
all such $z$.  For the homomorphism equation, use the first displayed
commutation equation and then the homomorphism field of $H$:
\[
 R^{\complement}(J(x),J(\mathsf{RightComp}(g)(x)))
 =R^{\complement}(H(K(x)),H(\mathsf{ShiftMap}(K(x)))) .
\]
Thus $J$ is a complement homomorphism for $\mathsf{RightComp}(g)$.  The
same argument with $L$ in place of $K$ transports a complement homomorphism
in the reverse direction.  Therefore the absence of the complement
homomorphism in the second clause of \noderef{GBetterRel} is equivalent to
the absence in \noderef{BetterRel}; the source actions are precisely those
of \noderef{RightComp} and \noderef{ShiftMap}.

It remains to justify, in the detail needed for the first clause of
\noderef{GBetterRel}, why absence of an $R$-homomorphic restriction is the
same as the existence of a complement homomorphism.  Fix a locally constant
$\varphi:\mathsf{IncSeq}\to A$ and define the Boolean coloring
\[
c(y)=
\begin{cases}
1&\text{if }R(\varphi(y),\varphi(\mathsf{RightComp}(g)(y))),\\
0&\text{otherwise.}
\end{cases}
\]
This is locally constant.  To check this at $y$, use
\noderef{LocallyConstant} to choose $n$ deciding $\varphi$ at $y$ and $m$
deciding $\varphi$ at $\mathsf{RightComp}(g)(y)$.  If $y'$ agrees with $y$
on the first $n$ entries, then $\varphi(y')=\varphi(y)$.  If it agrees on
the first $g(m)$ entries, then for each $i<m$ the order-embedding property
of $g$, recorded by \noderef{IncSeq}, gives $g(i)<g(m)$, and hence
\[
\mathsf{RightComp}(g)(y')(i)=y'(g(i))=y(g(i))
 =\mathsf{RightComp}(g)(y)(i).
\]
The evaluation identity uses \noderef{RightComp}.  Thus the two
right-composed sequences agree through $m$, so their
$\varphi$-values are equal as well.  By \noderef{PrefixAgree}, taking a
prefix long enough to satisfy both bounds makes both arguments of $R$ and
therefore $c$ constant.  This proves local constancy of $c$.

Let $F=F^c$ be the set of minimal deciding prefixes.  By
\noderef{DecidingFront}, $F$ is a front on $\mathbb N$.  By
\noderef{DecidingValue}, there is a map
$v:F\to\{0,1\}$ such that whenever $s\in F$ and
$\mathsf{ProperPrefixSet}(s,\operatorname{range}(x))$ in the sense of
\noderef{ProperPrefixSet},
\[
v(s)=c(x). \tag{1}
\]
Let $S=\{s\in F\mid v(s)=1\}$.  Applying \noderef{NashWilliams} to
$F$, its infinite base $\mathbb N$, and $S\subseteq F$ gives a front $F'$ on
an infinite set $Y$, with $F'\subseteq F$, such that either
$F'\subseteq S$ or $F'\cap S=\varnothing$.  The front conclusion and its
infinite-base clause are part of \noderef{Front}.  Choose, using
\noderef{InfiniteSetEnumeration}, an increasing enumeration $z$ of $Y$:
\[
\operatorname{range}(z)=Y. \tag{2}
\]

Fix $x\in\mathsf{IncSeq}$.  The range of
$\mathsf{IncSeqComp}(z,x)$ is an infinite subset of $Y$: it is the range of
an increasing enumeration by \noderef{IncSeq} and \noderef{IncSeqComp}, and
it is contained in $\operatorname{range}(z)=Y$.  The density clause of the
front \noderef{Front} therefore gives an $s_x\in F'$ with
\[
\mathsf{ProperPrefixSet}(s_x,
  \operatorname{range}(\mathsf{IncSeqComp}(z,x))). \tag{3}
\]
Since $s_x\in F'\subseteq F$, equation (1) gives
\[
v(s_x)=c(\mathsf{IncSeqComp}(z,x)). \tag{4}
\]
If $F'\subseteq S$, then (4) is $c(\mathsf{IncSeqComp}(z,x))=1$.
If $F'\cap S=\varnothing$, then $s_x\notin S$; because $s_x\in F$ and
$v(s_x)\in\{0,1\}$, (4) is instead
$c(\mathsf{IncSeqComp}(z,x))=0$.  Thus the coloring is constant on the
whole right ideal $z\circ\mathsf{IncSeq}$.

Put
\[
\psi=\mathsf{MultiSequenceRestrict}(\varphi,z),
\qquad
\psi(x)=\varphi(\mathsf{IncSeqComp}(z,x)).
\]
By \noderef{MultiSequenceRestrict} this is exactly restriction to the base
enumerated by $z$, and \noderef{RestrictionLocallyConstant} makes $\psi$
locally constant.  The required action identity is, for every $x$,
\[
\mathsf{RightComp}(g)(\mathsf{IncSeqComp}(z,x))
 =\mathsf{IncSeqComp}(z,\mathsf{RightComp}(g)(x)). \tag{5}
\]
Indeed, by \noderef{RightComp} and \noderef{IncSeqComp}, the two sides are
respectively $(z\circ x)\circ g$ and $z\circ(x\circ g)$, and these are
equal by associativity of composition.

If the homogeneous value above is $1$, the definition of $c$, equation
(5), and the displayed formula for $\psi$ give, for every $x$,
\[
R(\psi(x),\psi(\mathsf{RightComp}(g)(x))).
\]
Together with the local-constancy field already proved, these are exactly
the two fields of a \noderef{ContinuousRelHom}; taking $f=z$ supplies the
first clause of \noderef{GBetterRel} for $\varphi$.  If the homogeneous value
is $0$, the same calculation gives
\[
\neg R(\psi(x),\psi(\mathsf{RightComp}(g)(x)))
\quad\text{for every }x,
\]
so the same two fields instead give a continuous homomorphism into
$R^{\complement}$.  Consequently, if the second clause of
\noderef{GBetterRel} holds, the zero case is impossible and the first clause
holds for every locally constant $\varphi$.

Conversely, suppose that
$H:(\mathsf{IncSeq},\mathsf{RightComp}(g))\to(A,R^{\complement})$ is a
continuous homomorphism in the sense of \noderef{ContinuousRelHom}.  For
each $f\in\mathsf{IncSeq}$ define
\[
H_f=\mathsf{MultiSequenceRestrict}(H.\mathsf{toFun},f),
\qquad H_f(y)=H.\mathsf{toFun}(\mathsf{IncSeqComp}(f,y)).
\]
The local-constancy field of $H$ and
\noderef{RestrictionLocallyConstant} show that $H_f$ is locally constant.
The same associativity calculation as (5), using
\noderef{RightComp} and \noderef{IncSeqComp}, gives
\[
R^{\complement}(H_f(y),H_f(\mathsf{RightComp}(g)(y)))
\]
for every $y$, by the homomorphism field of $H$.  Hence every $H_f$ is a
continuous complement homomorphism.  If the first clause of
\noderef{GBetterRel} held for the locally constant map $H.\mathsf{toFun}$,
it would provide an $f$ and an $R$-homomorphism whose function is exactly
$H_f$.  At any input its homomorphism field would assert $R(a,b)$ while the
preceding paragraph asserts $R^{\complement}(a,b)=\neg R(a,b)$, a
contradiction.  Thus the first clause fails whenever a complement
homomorphism exists, and the two clauses in \noderef{GBetterRel} are
equivalent.  Combining this with the transport of the second clause to the
no-complement formulation in \noderef{BetterRel} proves the first displayed
equivalence.

Finally let $r$ be a quasi-order on $Q$.  By \noderef{BetterRel}, the
ordinary better-relation condition is the nonexistence of a
\noderef{ContinuousRelHom} from the shift into the complement relation.
Unpacking such a structure gives a locally constant map
$h:\mathsf{IncSeq}\to Q$ and, for every $x$,
\[
 \neg r(h(x),h(\mathsf{ShiftMap}(x))).
\]
By \noderef{BadMultiSequence}, this last condition says exactly that $h$
is bad for $r$.  Conversely, a locally constant bad multi-sequence supplies
the two fields of a \noderef{ContinuousRelHom} into the complement
relation.  Therefore the absence of this homomorphism is equivalent to the
statement in \noderef{Bqo} that every locally constant $h$ is not bad.
Thus $\operatorname{BetterRel}(r)\leftrightarrow\operatorname{Bqo}(r)$,
and the already proved general equivalence yields
\[
 \operatorname{GBetterRel}(g,r)\leftrightarrow\operatorname{Bqo}(r).
\]
\end{proof}
